% =============================================================================
% Chapter 1 — Background and Context
% Include from main.tex via:  % =============================================================================
% Chapter 1 — Background and Context
% Include from main.tex via:  % =============================================================================
% Chapter 1 — Background and Context
% Include from main.tex via:  % =============================================================================
% Chapter 1 — Background and Context
% Include from main.tex via:  \input{chap1}
% Do NOT compile this file alone. Set main.tex as the Overleaf main document.
% Required packages in main.tex: booktabs, graphicx, array
% Required in preamble: \graphicspath{{firas/}{figures/}}
% =============================================================================

\chapter{Background and Context}
\label{ch:background}

On many farms today, automation still works in simple chains. A sensor reads
soil moisture, a controller checks a threshold, and a pump switches on. Each
step makes sense on its own, but the loop rarely looks beyond the immediate
reading. It does not account for tomorrow's weather, the crop's growth stage, or
the longer-term trade-off between saving water and protecting yield.

That limitation becomes more visible as farms adopt more technologies. A single
field may now depend on soil sensors, weather data, communication networks,
simulation tools, and automated irrigation hardware, often supplied by different
vendors and updated on different schedules. The challenge is no longer to operate
one device well, but to make many independent parts work toward the same goal:
producing food efficiently while using water and nutrients responsibly. This is
where the \emph{System-of-Systems} (SoS) perspective becomes useful.

Figure~\ref{fig:traditional_vs_sos} illustrates the difference. The upper part
of the figure shows the traditional pattern: one input, one decision, one action.
The lower part shows a smart-farming SoS, in which weather data, field
observations, and crop-growth prediction feed a shared decision layer that
commands actuators and then receives updated feedback from the field. The rest
of this chapter builds the background needed for that view: SoS theory, smart
farming as an acknowledged SoS, crop simulation with DSSAT, and the reason
DSSAT is central to the work presented in this thesis.

\begin{figure}[ht]
  \centering
  \includegraphics[width=0.92\textwidth]{firas/traditional_vs_sos_farming.png}
  \caption{Comparison between traditional reactive farm automation (linear
    sensing--control--actuation chain) and a System-of-Systems smart farm, where
    weather, sensors, and crop simulation feed a shared decision layer that
    commands actuators in a closed loop.}
  \label{fig:traditional_vs_sos}
\end{figure}


\section{System of Systems}
\label{sec:sos}

A System of Systems (SoS) is not simply a large system. It is a set of
independent systems that keep their own purpose and operation, yet work together
to achieve something broader~\cite{maier1998architecting, jamshidi2008systems}.
A weather service, a sensor network, and a crop model can each function without
the others, but when they are combined under a common management goal, they can
support decisions that none of them could support alone~\cite{madni2014integration}.

Agriculture offers a clear example. A farmer may use soil probes from one
company, weather data from a public API, a crop model developed by a research
community, and irrigation equipment installed by a local provider. These tools
were not designed as one package, yet the farmer still needs them to fit
together. If they do not, management becomes fragmented: irrigation reacts to
moisture, fertilisation follows a calendar, and yield expectations ignore the
season's actual weather~\cite{nielsen2015sos, krisnawijaya2024architecture}.


\subsection{SoS Characteristics}

Maier~\cite{maier1998architecting} proposed five characteristics that help
distinguish an SoS from an ordinary integrated system. Table~\ref{tab:sos_characteristics}
lists them and translates each one into an agricultural setting.

\begin{table}[ht]
  \centering
  \caption{System-of-Systems characteristics and their meaning in smart farming.}
  \label{tab:sos_characteristics}
  \begin{tabular}{p{3.2cm} p{4.3cm} p{6.3cm}}
    \toprule
    \textbf{Characteristic} & \textbf{General meaning} & \textbf{Example in agriculture} \\
    \midrule
    Operational independence &
    Each system can function on its own &
    A weather API remains available even if irrigation valves are offline \\[4pt]

    Managerial independence &
    Different authorities own and manage subsystems &
    Sensors, DSSAT, and actuators are developed and maintained by different organisations \\[4pt]

    Geographic distribution &
    Systems are spread across locations &
    Field sensors, cloud services, and farmer terminals are spatially separated \\[4pt]

    Evolutionary development &
    Subsystems can be added, removed, or upgraded independently &
    A drone module or new crop model can be integrated without rebuilding the entire platform \\[4pt]

    Emergent behavior &
    Global behaviour exceeds local subsystem design &
    Yield and resource efficiency emerge from the interaction of weather, soil, crop, and management decisions \\
    \bottomrule
  \end{tabular}
\end{table}

\paragraph{Operational independence.}
Each part of an SoS should remain useful on its own. A weather station keeps
recording even when irrigation is switched off. A crop model can still project
yield from historical weather even if live sensors fail. That independence matters
in agriculture because farms rarely replace all of their equipment at once.

\paragraph{Managerial independence.}
In practice, no single organisation controls every layer of a smart farm. Sensor
manufacturers, weather providers, software platforms, and public crop models
such as DSSAT follow their own development paths. Integration therefore depends
on interfaces and standards, not on one central redesign of the whole
stack~\cite{madni2014integration}.

\paragraph{Geographic distribution.}
Field sensors sit in the soil, gateways may be located at the farm edge, cloud
services run elsewhere, and the farmer may check results on a phone. Coordination
has to happen across distance and networks, not inside one local
machine~\cite{nielsen2015sos}.

\paragraph{Evolutionary development.}
Farm technologies do not age at the same speed. A farmer may upgrade soil probes
one year, adopt a new forecast service the next, and keep the same irrigation
hardware for a decade. An SoS architecture needs to absorb those changes without
forcing a complete restart~\cite{nielsen2015sos, madni2014integration}.

\paragraph{Emergent behavior.}
Yield, water productivity, and nitrogen efficiency are not produced by any single
device. They emerge from the interaction between weather, soil conditions, crop
development, and management over time~\cite{jamshidi2008systems, boote2019cropmodelling}.
That is also why agricultural decision-making is difficult: the outcome farmers
care about appears at the system level, not inside one subsystem.


\subsection{Classification of SoS}
\label{subsec:sos_classification}

Not all Systems of Systems are organised in the same way. Maier~\cite{maier1998architecting}
classified them according to how much central authority exists over the
constituent systems. He originally identified three types---directed,
collaborative, and virtual---and a fourth type, \emph{acknowledged}, was added
later by Dahmann and Baldwin and incorporated into ISO~21839~\cite{nielsen2015sos}.
Table~\ref{tab:sos_taxonomy} gives a short overview.

\begin{table}[ht]
  \centering
  \caption{Classification of Systems of Systems (Maier; extended by Dahmann and Baldwin).}
  \label{tab:sos_taxonomy}
  \begin{tabular}{p{2.5cm} p{4.5cm} p{6.5cm}}
    \toprule
    \textbf{Type} & \textbf{Management structure} & \textbf{Typical agricultural interpretation} \\
    \midrule
    Directed &
    A central authority creates and manages the SoS &
    A fully integrated farm platform owned and operated by one company \\[4pt]

    Acknowledged &
    A recognised SoS manager coordinates independently owned systems &
    A farm-management platform integrates vendor sensors, weather data, DSSAT, and actuators through agreed interfaces \\[4pt]

    Collaborative &
    No central manager; systems cooperate voluntarily &
    A farmer manually combines advisory tools, weather services, and field monitoring applications \\[4pt]

    Virtual &
    No manager and no shared purpose &
    Independent technologies used side by side without real integration \\
    \bottomrule
  \end{tabular}
\end{table}

In a directed SoS, the central authority defines the mission and the subsystems
are built to serve it. In an acknowledged SoS, a coordinating layer exists, but
the subsystems still belong to different owners and evolve on their own;
integration happens through agreement rather than command. Collaborative SoS
depend on voluntary cooperation, while virtual SoS show emergent behaviour even
when the parts were never intended to form a unified
system~\cite{maier1998architecting, nielsen2015sos}.

For modern smart-farming platforms, the acknowledged type is especially
convincing. A farm dashboard, digital twin, or decision-support service can
coordinate tools it does not own. That description fits agriculture well: the
farmer or platform operator manages the overall mission, but the sensors,
models, and actuators remain separate products developed by separate actors.


\section{Smart Farming as an Acknowledged SoS}
\label{sec:smart_farming_sos}

Precision agriculture began with a simple idea: fields are not uniform, so
management should not be uniform either~\cite{aarzoo2025precision, kushwaha2024precision}.
Instead of applying the same irrigation or fertiliser rate everywhere, farmers
can adjust treatment to local soil conditions, moisture levels, crop status, and
weather~\cite{kushwaha2024precision}. Over time, that idea grew into smart
farming: a setting in which sensing, connectivity, data analysis, simulation,
and automated actuation are woven into one management
environment~\cite{hayat2025iot, jadhav2023smartfarming}.

GPS, GIS, remote sensing, and IoT platforms now make it possible to observe
field variability with much greater detail than before~\cite{sharma2020ml,
adewuyi2024datascience}. The result is no longer a single machine or a single
software package, but a network of tools that must cooperate. That is why smart
farming is better described as an SoS than as a standalone automation
system. Table~\ref{tab:traditional_vs_sos} contrasts the older reactive pattern
with this integrated view.

\begin{table}[ht]
  \centering
  \caption{Comparison between traditional farm automation and System-of-Systems smart farming.}
  \label{tab:traditional_vs_sos}
  \begin{tabular}{p{3.5cm} p{5.3cm} p{5.3cm}}
    \toprule
    \textbf{Dimension} & \textbf{Traditional approach} & \textbf{SoS smart farming} \\
    \midrule
    Architecture &
    Isolated chain: sensor $\rightarrow$ controller $\rightarrow$ actuator &
    Multiple independent subsystems connected through a shared decision layer \\[4pt]

    Decision logic &
    Reactive rule-based control &
    Predictive and integrative decision-making \\[4pt]

    Information flow &
    One input, one action &
    Weather, field observations, and crop-model outputs are fused \\[4pt]

    Time horizon &
    Immediate response &
    Short-term control combined with seasonal crop response \\[4pt]

    Subsystem autonomy &
    Low; components designed as one package &
    High; each subsystem can exist independently \\[4pt]

    Scalability &
    Limited to one loop or one field function &
    New sensors, models, or actuators can be added incrementally \\[4pt]

    Closed loop &
    Local physical feedback only &
    Physical feedback plus model-based anticipation \\
    \bottomrule
  \end{tabular}
\end{table}


\subsection{Constituent Systems in Smart Farming}

Even when they are not sold as separate products, the main parts of a smart farm
play distinct roles. Table~\ref{tab:smart_farm_subsystems} summarises the
subsystems that appear repeatedly in the literature and in practical deployments.

\begin{table}[ht]
  \centering
  \caption{Main constituent systems in a smart farming System-of-Systems.}
  \label{tab:smart_farm_subsystems}
  \begin{tabular}{p{2.8cm} p{3.5cm} p{3.5cm} p{3.8cm}}
    \toprule
    \textbf{Subsystem} & \textbf{Main role} & \textbf{Typical inputs} & \textbf{Typical outputs} \\
    \midrule
    Weather system &
    Climate forcing and forecasting &
    Historical weather, numerical forecasts &
    Temperature, rainfall, solar radiation \\[4pt]

    Sensor network &
    Field observation &
    Soil moisture, temperature, nutrient status &
    Ground-truth state measurements \\[4pt]

    Crop simulation model &
    Growth and stress prediction &
    Weather, soil, cultivar, management actions &
    Crop condition, stress indicators, yield projections \\[4pt]

    Decision layer &
    Policy optimisation &
    Observations, model outputs, objectives &
    Irrigation and fertilisation decisions \\[4pt]

    Actuators &
    Execution of management actions &
    Control commands &
    Applied water and nutrients \\
    \bottomrule
  \end{tabular}
\end{table}

\paragraph{Weather system.}
Weather remains the main external driver of crop growth. Local stations can
supply temperature, rainfall, and radiation, while services such as the NASA
POWER API can fill gaps with historical or near-real-time
data~\cite{hayat2025iot, codeluppi2020lorafarm}. Without reliable weather input,
both irrigation planning and crop modelling become unreliable.

\paragraph{Sensor network.}
Soil sensors based on capacitance or time-domain reflectometry (TDR) can track
moisture, temperature, and nitrogen status at different
depths~\cite{hayat2025iot, jadhav2023smartfarming}. They provide the field-level
observations that make site-specific management possible.

\paragraph{Communication infrastructure.}
Distributed farms need a way to move data from the field to the decision layer.
LPWAN technologies, especially LoRaWAN, are often used because they cover long
distances with low energy use~\cite{codeluppi2020lorafarm, jadhav2023smartfarming}.
In an SoS view, this communication layer is what turns separate devices into a
connected system.

\paragraph{Crop simulation model.}
Sensors show what is happening now; crop models help estimate what management
will mean later in the season. They connect weather, soil, and farmer actions to
likely crop development, stress, and harvest
outcomes~\cite{boote2019cropmodelling, hoogenboom2019dssat}. They extend the
decision horizon beyond the latest sensor reading.

\paragraph{Decision layer.}
At some point, the farm must choose an action: irrigate or wait, fertilise or
hold. That choice may come from a rule engine, an optimisation tool, a human
adviser, or---as in this thesis---a learning-based controller. Whatever form it
takes, the decision layer is where the SoS becomes operational.

\paragraph{Actuators.}
Irrigation valves and variable-rate fertiliser applicators turn decisions into
physical changes in the field~\cite{hayat2025iot}. Once water or nitrogen is
applied, the crop and soil state changes, and the loop starts again.


\subsection{Why Smart Farming Fits the Acknowledged SoS Model}

Smart farming resembles an acknowledged SoS for a simple reason: coordination
exists, but ownership stays divided. A platform or farm manager integrates
weather, sensors, models, and actuators, yet those components usually come from
different suppliers and remain replaceable. A farmer can change sensor brand or
forecast provider without rebuilding the entire system.

That does not mean every farm already works this way. Many operations still look
more like the traditional chain in Figure~\ref{fig:traditional_vs_sos}. The
acknowledged SoS model is therefore best seen as the direction in which mature
smart-farming platforms are heading: separate tools, shared
mission~\cite{krisnawijaya2024architecture, nielsen2015sos}.

There is another practical complication. The data arriving from these subsystems
do not share the same form. Sensor streams are time series, drone imagery is
visual and high-dimensional, and crop-model outputs may appear as state
variables or structured reports~\cite{krisnawijaya2024architecture,
adewuyi2024datascience}. A workable SoS therefore needs both communication
infrastructure and a decision layer able to bring those different sources
together.


\section{Crop Simulation and DSSAT}
\label{sec:dssat}

Sensors alone cannot answer every management question. A moisture reading may
show that the soil is dry today, but it does not by itself predict final yield,
nitrogen losses, or crop stress two weeks later. Crop simulation models fill
that gap by representing how plants, soil water, and nutrients respond over
time~\cite{boote2019cropmodelling}.

Among the available platforms, the Decision Support System for Agrotechnology
Transfer (DSSAT) is one of the most established. It has been used for decades to
combine variety, environment, and management information in dynamic crop
simulations~\cite{hoogenboom2019dssat, sachin2019dssat}. Its main value is not
a single yield number at harvest, but the ability to follow how daily weather
and management actions shape crop response over the season.


\subsection{The G$\times$E$\times$M Framework}

DSSAT is built around the Genotype--Environment--Management (G$\times$E$\times$M)
framework. Crop performance depends on all three. A high-yielding cultivar still
needs suitable weather and reasonable management; good irrigation cannot fully
compensate for a poor variety or badly timed
operations~\cite{boote2019cropmodelling}. Table~\ref{tab:gxem} breaks down the
three components.

\begin{table}[ht]
  \centering
  \caption{The Genotype--Environment--Management (G$\times$E$\times$M) framework in DSSAT.}
  \label{tab:gxem}
  \begin{tabular}{p{2.5cm} p{4.3cm} p{6.5cm}}
    \toprule
    \textbf{Factor} & \textbf{Meaning} & \textbf{Examples in DSSAT} \\
    \midrule
    Genotype (G) &
    Crop genetic potential and development pattern &
    Cultivar parameters, phenology, grain-fill capacity \\[4pt]

    Environment (E) &
    Climatic and edaphic conditions &
    Daily weather, soil layers, initial water and nitrogen status \\[4pt]

    Management (M) &
    Farmer interventions during the season &
    Planting date, irrigation depth, nitrogen applications \\
    \bottomrule
  \end{tabular}
\end{table}

This matters for decision support because management actions do not influence
yield directly. Irrigation changes soil moisture; fertilisation changes nutrient
availability; those changes then affect how the crop grows, whether it becomes
stressed, and what is harvested at the end of the season. DSSAT makes that chain
explicit.


\subsection{The DSSAT Cropping System Model}

At the heart of DSSAT is the Cropping System Model, usually referred to as CSM.
Instead of jumping straight to a harvest number, CSM moves through the season one
day at a time. On each day, it reads the weather, updates soil moisture and
nutrient conditions, applies any field operations such as irrigation or
fertiliser, and then updates the state of the
crop~\cite{hoogenboom2019dssat}.

For the maize experiments in this thesis, the relevant crop model is
CERES-Maize. In everyday terms, it tracks how the crop develops after planting,
how well it copes with dry or nutrient-poor conditions, and how much grain is
likely to be produced by harvest. That makes it useful for the kind of management
questions studied here, where water and nitrogen decisions need to be judged not
only for their immediate effect, but also for their impact over the rest of the
season~\cite{boote2019cropmodelling, sachin2019dssat}.

This is also what distinguishes DSSAT from a simple prediction formula. A basic
model might estimate yield from a few inputs, but it would miss the interactions
that make farm management difficult in practice: too little water early in the
season, too much nitrogen before rain, or the opposite. CSM is built to represent
those interactions rather than ignore them.

In most published DSSAT studies, the management plan is decided in advance and the
model is run as a full-season experiment. In this thesis, the same crop-response
logic is used, but irrigation and nitrogen applications may change over time. The
software details of that setup are presented in the methodology chapters; here,
the important point is simply that CSM provides the crop-response core around
which the rest of the system is built.


\section{Why DSSAT for This Work}
\label{sec:dssat_architecture}

Look again at Figure~\ref{fig:traditional_vs_sos}. Weather tells the system what
the climate is doing. Sensors describe the current field state. DSSAT connects
those conditions to likely crop response under different irrigation and
fertilisation choices. It is the part of the SoS that translates management
decisions into crop outcomes.

That is the question at the centre of this thesis. The work is not about
building a complete smart farm with its own communication stack, drone service,
and farm dashboard. It is about what happens inside the SoS when yield, water
use, and nitrogen efficiency cannot all be maximised at once: how should
irrigation and fertilisation be adjusted over the season? DSSAT makes that
question concrete because it allows those trade-offs to be explored against a
realistic crop-response model~\cite{hoogenboom2019dssat, sachin2019dssat}.

The next chapter introduces the reinforcement-learning framework used to study
that problem. The methodology chapters then explain how daily management
decisions are exchanged with the simulator in practice.


\section{Conclusion}
\label{sec:ch1_conclusion}

Modern farming is moving away from isolated devices that react to a single
reading and toward environments in which weather services, sensors, crop models,
and actuators have to work together. This chapter used the System-of-Systems view
to describe that shift and to explain why coordination alone is not enough:
some part of the system still has to decide how water, nitrogen, and crop
outcomes should be balanced over a season.

Smart farming fits the acknowledged SoS model because farms increasingly rely on
tools they do not fully own, yet still need those tools to serve a common
management goal. DSSAT enters the picture as the crop-simulation part of that
picture. It does not replace sensors or actuators, but it helps answer what
happens to the crop when management changes under real weather and soil
conditions.

That is why DSSAT matters for this thesis. The work is less about assembling a
complete commercial smart farm and more about studying a difficult decision
problem inside that broader setting: how irrigation and fertilisation should
adapt when yield, water use, and nitrogen efficiency pull in different
directions. The next chapter takes up that decision problem directly and
introduces the reinforcement-learning framework used to address it.

% Do NOT compile this file alone. Set main.tex as the Overleaf main document.
% Required packages in main.tex: booktabs, graphicx, array
% Required in preamble: \graphicspath{{firas/}{figures/}}
% =============================================================================

\chapter{Background and Context}
\label{ch:background}

On many farms today, automation still works in simple chains. A sensor reads
soil moisture, a controller checks a threshold, and a pump switches on. Each
step makes sense on its own, but the loop rarely looks beyond the immediate
reading. It does not account for tomorrow's weather, the crop's growth stage, or
the longer-term trade-off between saving water and protecting yield.

That limitation becomes more visible as farms adopt more technologies. A single
field may now depend on soil sensors, weather data, communication networks,
simulation tools, and automated irrigation hardware, often supplied by different
vendors and updated on different schedules. The challenge is no longer to operate
one device well, but to make many independent parts work toward the same goal:
producing food efficiently while using water and nutrients responsibly. This is
where the \emph{System-of-Systems} (SoS) perspective becomes useful.

Figure~\ref{fig:traditional_vs_sos} illustrates the difference. The upper part
of the figure shows the traditional pattern: one input, one decision, one action.
The lower part shows a smart-farming SoS, in which weather data, field
observations, and crop-growth prediction feed a shared decision layer that
commands actuators and then receives updated feedback from the field. The rest
of this chapter builds the background needed for that view: SoS theory, smart
farming as an acknowledged SoS, crop simulation with DSSAT, and the reason
DSSAT is central to the work presented in this thesis.

\begin{figure}[ht]
  \centering
  \includegraphics[width=0.92\textwidth]{firas/traditional_vs_sos_farming.png}
  \caption{Comparison between traditional reactive farm automation (linear
    sensing--control--actuation chain) and a System-of-Systems smart farm, where
    weather, sensors, and crop simulation feed a shared decision layer that
    commands actuators in a closed loop.}
  \label{fig:traditional_vs_sos}
\end{figure}


\section{System of Systems}
\label{sec:sos}

A System of Systems (SoS) is not simply a large system. It is a set of
independent systems that keep their own purpose and operation, yet work together
to achieve something broader~\cite{maier1998architecting, jamshidi2008systems}.
A weather service, a sensor network, and a crop model can each function without
the others, but when they are combined under a common management goal, they can
support decisions that none of them could support alone~\cite{madni2014integration}.

Agriculture offers a clear example. A farmer may use soil probes from one
company, weather data from a public API, a crop model developed by a research
community, and irrigation equipment installed by a local provider. These tools
were not designed as one package, yet the farmer still needs them to fit
together. If they do not, management becomes fragmented: irrigation reacts to
moisture, fertilisation follows a calendar, and yield expectations ignore the
season's actual weather~\cite{nielsen2015sos, krisnawijaya2024architecture}.


\subsection{SoS Characteristics}

Maier~\cite{maier1998architecting} proposed five characteristics that help
distinguish an SoS from an ordinary integrated system. Table~\ref{tab:sos_characteristics}
lists them and translates each one into an agricultural setting.

\begin{table}[ht]
  \centering
  \caption{System-of-Systems characteristics and their meaning in smart farming.}
  \label{tab:sos_characteristics}
  \begin{tabular}{p{3.2cm} p{4.3cm} p{6.3cm}}
    \toprule
    \textbf{Characteristic} & \textbf{General meaning} & \textbf{Example in agriculture} \\
    \midrule
    Operational independence &
    Each system can function on its own &
    A weather API remains available even if irrigation valves are offline \\[4pt]

    Managerial independence &
    Different authorities own and manage subsystems &
    Sensors, DSSAT, and actuators are developed and maintained by different organisations \\[4pt]

    Geographic distribution &
    Systems are spread across locations &
    Field sensors, cloud services, and farmer terminals are spatially separated \\[4pt]

    Evolutionary development &
    Subsystems can be added, removed, or upgraded independently &
    A drone module or new crop model can be integrated without rebuilding the entire platform \\[4pt]

    Emergent behavior &
    Global behaviour exceeds local subsystem design &
    Yield and resource efficiency emerge from the interaction of weather, soil, crop, and management decisions \\
    \bottomrule
  \end{tabular}
\end{table}

\paragraph{Operational independence.}
Each part of an SoS should remain useful on its own. A weather station keeps
recording even when irrigation is switched off. A crop model can still project
yield from historical weather even if live sensors fail. That independence matters
in agriculture because farms rarely replace all of their equipment at once.

\paragraph{Managerial independence.}
In practice, no single organisation controls every layer of a smart farm. Sensor
manufacturers, weather providers, software platforms, and public crop models
such as DSSAT follow their own development paths. Integration therefore depends
on interfaces and standards, not on one central redesign of the whole
stack~\cite{madni2014integration}.

\paragraph{Geographic distribution.}
Field sensors sit in the soil, gateways may be located at the farm edge, cloud
services run elsewhere, and the farmer may check results on a phone. Coordination
has to happen across distance and networks, not inside one local
machine~\cite{nielsen2015sos}.

\paragraph{Evolutionary development.}
Farm technologies do not age at the same speed. A farmer may upgrade soil probes
one year, adopt a new forecast service the next, and keep the same irrigation
hardware for a decade. An SoS architecture needs to absorb those changes without
forcing a complete restart~\cite{nielsen2015sos, madni2014integration}.

\paragraph{Emergent behavior.}
Yield, water productivity, and nitrogen efficiency are not produced by any single
device. They emerge from the interaction between weather, soil conditions, crop
development, and management over time~\cite{jamshidi2008systems, boote2019cropmodelling}.
That is also why agricultural decision-making is difficult: the outcome farmers
care about appears at the system level, not inside one subsystem.


\subsection{Classification of SoS}
\label{subsec:sos_classification}

Not all Systems of Systems are organised in the same way. Maier~\cite{maier1998architecting}
classified them according to how much central authority exists over the
constituent systems. He originally identified three types---directed,
collaborative, and virtual---and a fourth type, \emph{acknowledged}, was added
later by Dahmann and Baldwin and incorporated into ISO~21839~\cite{nielsen2015sos}.
Table~\ref{tab:sos_taxonomy} gives a short overview.

\begin{table}[ht]
  \centering
  \caption{Classification of Systems of Systems (Maier; extended by Dahmann and Baldwin).}
  \label{tab:sos_taxonomy}
  \begin{tabular}{p{2.5cm} p{4.5cm} p{6.5cm}}
    \toprule
    \textbf{Type} & \textbf{Management structure} & \textbf{Typical agricultural interpretation} \\
    \midrule
    Directed &
    A central authority creates and manages the SoS &
    A fully integrated farm platform owned and operated by one company \\[4pt]

    Acknowledged &
    A recognised SoS manager coordinates independently owned systems &
    A farm-management platform integrates vendor sensors, weather data, DSSAT, and actuators through agreed interfaces \\[4pt]

    Collaborative &
    No central manager; systems cooperate voluntarily &
    A farmer manually combines advisory tools, weather services, and field monitoring applications \\[4pt]

    Virtual &
    No manager and no shared purpose &
    Independent technologies used side by side without real integration \\
    \bottomrule
  \end{tabular}
\end{table}

In a directed SoS, the central authority defines the mission and the subsystems
are built to serve it. In an acknowledged SoS, a coordinating layer exists, but
the subsystems still belong to different owners and evolve on their own;
integration happens through agreement rather than command. Collaborative SoS
depend on voluntary cooperation, while virtual SoS show emergent behaviour even
when the parts were never intended to form a unified
system~\cite{maier1998architecting, nielsen2015sos}.

For modern smart-farming platforms, the acknowledged type is especially
convincing. A farm dashboard, digital twin, or decision-support service can
coordinate tools it does not own. That description fits agriculture well: the
farmer or platform operator manages the overall mission, but the sensors,
models, and actuators remain separate products developed by separate actors.


\section{Smart Farming as an Acknowledged SoS}
\label{sec:smart_farming_sos}

Precision agriculture began with a simple idea: fields are not uniform, so
management should not be uniform either~\cite{aarzoo2025precision, kushwaha2024precision}.
Instead of applying the same irrigation or fertiliser rate everywhere, farmers
can adjust treatment to local soil conditions, moisture levels, crop status, and
weather~\cite{kushwaha2024precision}. Over time, that idea grew into smart
farming: a setting in which sensing, connectivity, data analysis, simulation,
and automated actuation are woven into one management
environment~\cite{hayat2025iot, jadhav2023smartfarming}.

GPS, GIS, remote sensing, and IoT platforms now make it possible to observe
field variability with much greater detail than before~\cite{sharma2020ml,
adewuyi2024datascience}. The result is no longer a single machine or a single
software package, but a network of tools that must cooperate. That is why smart
farming is better described as an SoS than as a standalone automation
system. Table~\ref{tab:traditional_vs_sos} contrasts the older reactive pattern
with this integrated view.

\begin{table}[ht]
  \centering
  \caption{Comparison between traditional farm automation and System-of-Systems smart farming.}
  \label{tab:traditional_vs_sos}
  \begin{tabular}{p{3.5cm} p{5.3cm} p{5.3cm}}
    \toprule
    \textbf{Dimension} & \textbf{Traditional approach} & \textbf{SoS smart farming} \\
    \midrule
    Architecture &
    Isolated chain: sensor $\rightarrow$ controller $\rightarrow$ actuator &
    Multiple independent subsystems connected through a shared decision layer \\[4pt]

    Decision logic &
    Reactive rule-based control &
    Predictive and integrative decision-making \\[4pt]

    Information flow &
    One input, one action &
    Weather, field observations, and crop-model outputs are fused \\[4pt]

    Time horizon &
    Immediate response &
    Short-term control combined with seasonal crop response \\[4pt]

    Subsystem autonomy &
    Low; components designed as one package &
    High; each subsystem can exist independently \\[4pt]

    Scalability &
    Limited to one loop or one field function &
    New sensors, models, or actuators can be added incrementally \\[4pt]

    Closed loop &
    Local physical feedback only &
    Physical feedback plus model-based anticipation \\
    \bottomrule
  \end{tabular}
\end{table}


\subsection{Constituent Systems in Smart Farming}

Even when they are not sold as separate products, the main parts of a smart farm
play distinct roles. Table~\ref{tab:smart_farm_subsystems} summarises the
subsystems that appear repeatedly in the literature and in practical deployments.

\begin{table}[ht]
  \centering
  \caption{Main constituent systems in a smart farming System-of-Systems.}
  \label{tab:smart_farm_subsystems}
  \begin{tabular}{p{2.8cm} p{3.5cm} p{3.5cm} p{3.8cm}}
    \toprule
    \textbf{Subsystem} & \textbf{Main role} & \textbf{Typical inputs} & \textbf{Typical outputs} \\
    \midrule
    Weather system &
    Climate forcing and forecasting &
    Historical weather, numerical forecasts &
    Temperature, rainfall, solar radiation \\[4pt]

    Sensor network &
    Field observation &
    Soil moisture, temperature, nutrient status &
    Ground-truth state measurements \\[4pt]

    Crop simulation model &
    Growth and stress prediction &
    Weather, soil, cultivar, management actions &
    Crop condition, stress indicators, yield projections \\[4pt]

    Decision layer &
    Policy optimisation &
    Observations, model outputs, objectives &
    Irrigation and fertilisation decisions \\[4pt]

    Actuators &
    Execution of management actions &
    Control commands &
    Applied water and nutrients \\
    \bottomrule
  \end{tabular}
\end{table}

\paragraph{Weather system.}
Weather remains the main external driver of crop growth. Local stations can
supply temperature, rainfall, and radiation, while services such as the NASA
POWER API can fill gaps with historical or near-real-time
data~\cite{hayat2025iot, codeluppi2020lorafarm}. Without reliable weather input,
both irrigation planning and crop modelling become unreliable.

\paragraph{Sensor network.}
Soil sensors based on capacitance or time-domain reflectometry (TDR) can track
moisture, temperature, and nitrogen status at different
depths~\cite{hayat2025iot, jadhav2023smartfarming}. They provide the field-level
observations that make site-specific management possible.

\paragraph{Communication infrastructure.}
Distributed farms need a way to move data from the field to the decision layer.
LPWAN technologies, especially LoRaWAN, are often used because they cover long
distances with low energy use~\cite{codeluppi2020lorafarm, jadhav2023smartfarming}.
In an SoS view, this communication layer is what turns separate devices into a
connected system.

\paragraph{Crop simulation model.}
Sensors show what is happening now; crop models help estimate what management
will mean later in the season. They connect weather, soil, and farmer actions to
likely crop development, stress, and harvest
outcomes~\cite{boote2019cropmodelling, hoogenboom2019dssat}. They extend the
decision horizon beyond the latest sensor reading.

\paragraph{Decision layer.}
At some point, the farm must choose an action: irrigate or wait, fertilise or
hold. That choice may come from a rule engine, an optimisation tool, a human
adviser, or---as in this thesis---a learning-based controller. Whatever form it
takes, the decision layer is where the SoS becomes operational.

\paragraph{Actuators.}
Irrigation valves and variable-rate fertiliser applicators turn decisions into
physical changes in the field~\cite{hayat2025iot}. Once water or nitrogen is
applied, the crop and soil state changes, and the loop starts again.


\subsection{Why Smart Farming Fits the Acknowledged SoS Model}

Smart farming resembles an acknowledged SoS for a simple reason: coordination
exists, but ownership stays divided. A platform or farm manager integrates
weather, sensors, models, and actuators, yet those components usually come from
different suppliers and remain replaceable. A farmer can change sensor brand or
forecast provider without rebuilding the entire system.

That does not mean every farm already works this way. Many operations still look
more like the traditional chain in Figure~\ref{fig:traditional_vs_sos}. The
acknowledged SoS model is therefore best seen as the direction in which mature
smart-farming platforms are heading: separate tools, shared
mission~\cite{krisnawijaya2024architecture, nielsen2015sos}.

There is another practical complication. The data arriving from these subsystems
do not share the same form. Sensor streams are time series, drone imagery is
visual and high-dimensional, and crop-model outputs may appear as state
variables or structured reports~\cite{krisnawijaya2024architecture,
adewuyi2024datascience}. A workable SoS therefore needs both communication
infrastructure and a decision layer able to bring those different sources
together.


\section{Crop Simulation and DSSAT}
\label{sec:dssat}

Sensors alone cannot answer every management question. A moisture reading may
show that the soil is dry today, but it does not by itself predict final yield,
nitrogen losses, or crop stress two weeks later. Crop simulation models fill
that gap by representing how plants, soil water, and nutrients respond over
time~\cite{boote2019cropmodelling}.

Among the available platforms, the Decision Support System for Agrotechnology
Transfer (DSSAT) is one of the most established. It has been used for decades to
combine variety, environment, and management information in dynamic crop
simulations~\cite{hoogenboom2019dssat, sachin2019dssat}. Its main value is not
a single yield number at harvest, but the ability to follow how daily weather
and management actions shape crop response over the season.


\subsection{The G$\times$E$\times$M Framework}

DSSAT is built around the Genotype--Environment--Management (G$\times$E$\times$M)
framework. Crop performance depends on all three. A high-yielding cultivar still
needs suitable weather and reasonable management; good irrigation cannot fully
compensate for a poor variety or badly timed
operations~\cite{boote2019cropmodelling}. Table~\ref{tab:gxem} breaks down the
three components.

\begin{table}[ht]
  \centering
  \caption{The Genotype--Environment--Management (G$\times$E$\times$M) framework in DSSAT.}
  \label{tab:gxem}
  \begin{tabular}{p{2.5cm} p{4.3cm} p{6.5cm}}
    \toprule
    \textbf{Factor} & \textbf{Meaning} & \textbf{Examples in DSSAT} \\
    \midrule
    Genotype (G) &
    Crop genetic potential and development pattern &
    Cultivar parameters, phenology, grain-fill capacity \\[4pt]

    Environment (E) &
    Climatic and edaphic conditions &
    Daily weather, soil layers, initial water and nitrogen status \\[4pt]

    Management (M) &
    Farmer interventions during the season &
    Planting date, irrigation depth, nitrogen applications \\
    \bottomrule
  \end{tabular}
\end{table}

This matters for decision support because management actions do not influence
yield directly. Irrigation changes soil moisture; fertilisation changes nutrient
availability; those changes then affect how the crop grows, whether it becomes
stressed, and what is harvested at the end of the season. DSSAT makes that chain
explicit.


\subsection{The DSSAT Cropping System Model}

At the heart of DSSAT is the Cropping System Model, usually referred to as CSM.
Instead of jumping straight to a harvest number, CSM moves through the season one
day at a time. On each day, it reads the weather, updates soil moisture and
nutrient conditions, applies any field operations such as irrigation or
fertiliser, and then updates the state of the
crop~\cite{hoogenboom2019dssat}.

For the maize experiments in this thesis, the relevant crop model is
CERES-Maize. In everyday terms, it tracks how the crop develops after planting,
how well it copes with dry or nutrient-poor conditions, and how much grain is
likely to be produced by harvest. That makes it useful for the kind of management
questions studied here, where water and nitrogen decisions need to be judged not
only for their immediate effect, but also for their impact over the rest of the
season~\cite{boote2019cropmodelling, sachin2019dssat}.

This is also what distinguishes DSSAT from a simple prediction formula. A basic
model might estimate yield from a few inputs, but it would miss the interactions
that make farm management difficult in practice: too little water early in the
season, too much nitrogen before rain, or the opposite. CSM is built to represent
those interactions rather than ignore them.

In most published DSSAT studies, the management plan is decided in advance and the
model is run as a full-season experiment. In this thesis, the same crop-response
logic is used, but irrigation and nitrogen applications may change over time. The
software details of that setup are presented in the methodology chapters; here,
the important point is simply that CSM provides the crop-response core around
which the rest of the system is built.


\section{Why DSSAT for This Work}
\label{sec:dssat_architecture}

Look again at Figure~\ref{fig:traditional_vs_sos}. Weather tells the system what
the climate is doing. Sensors describe the current field state. DSSAT connects
those conditions to likely crop response under different irrigation and
fertilisation choices. It is the part of the SoS that translates management
decisions into crop outcomes.

That is the question at the centre of this thesis. The work is not about
building a complete smart farm with its own communication stack, drone service,
and farm dashboard. It is about what happens inside the SoS when yield, water
use, and nitrogen efficiency cannot all be maximised at once: how should
irrigation and fertilisation be adjusted over the season? DSSAT makes that
question concrete because it allows those trade-offs to be explored against a
realistic crop-response model~\cite{hoogenboom2019dssat, sachin2019dssat}.

The next chapter introduces the reinforcement-learning framework used to study
that problem. The methodology chapters then explain how daily management
decisions are exchanged with the simulator in practice.


\section{Conclusion}
\label{sec:ch1_conclusion}

Modern farming is moving away from isolated devices that react to a single
reading and toward environments in which weather services, sensors, crop models,
and actuators have to work together. This chapter used the System-of-Systems view
to describe that shift and to explain why coordination alone is not enough:
some part of the system still has to decide how water, nitrogen, and crop
outcomes should be balanced over a season.

Smart farming fits the acknowledged SoS model because farms increasingly rely on
tools they do not fully own, yet still need those tools to serve a common
management goal. DSSAT enters the picture as the crop-simulation part of that
picture. It does not replace sensors or actuators, but it helps answer what
happens to the crop when management changes under real weather and soil
conditions.

That is why DSSAT matters for this thesis. The work is less about assembling a
complete commercial smart farm and more about studying a difficult decision
problem inside that broader setting: how irrigation and fertilisation should
adapt when yield, water use, and nitrogen efficiency pull in different
directions. The next chapter takes up that decision problem directly and
introduces the reinforcement-learning framework used to address it.

% Do NOT compile this file alone. Set main.tex as the Overleaf main document.
% Required packages in main.tex: booktabs, graphicx, array
% Required in preamble: \graphicspath{{firas/}{figures/}}
% =============================================================================

\chapter{Background and Context}
\label{ch:background}

On many farms today, automation still works in simple chains. A sensor reads
soil moisture, a controller checks a threshold, and a pump switches on. Each
step makes sense on its own, but the loop rarely looks beyond the immediate
reading. It does not account for tomorrow's weather, the crop's growth stage, or
the longer-term trade-off between saving water and protecting yield.

That limitation becomes more visible as farms adopt more technologies. A single
field may now depend on soil sensors, weather data, communication networks,
simulation tools, and automated irrigation hardware, often supplied by different
vendors and updated on different schedules. The challenge is no longer to operate
one device well, but to make many independent parts work toward the same goal:
producing food efficiently while using water and nutrients responsibly. This is
where the \emph{System-of-Systems} (SoS) perspective becomes useful.

Figure~\ref{fig:traditional_vs_sos} illustrates the difference. The upper part
of the figure shows the traditional pattern: one input, one decision, one action.
The lower part shows a smart-farming SoS, in which weather data, field
observations, and crop-growth prediction feed a shared decision layer that
commands actuators and then receives updated feedback from the field. The rest
of this chapter builds the background needed for that view: SoS theory, smart
farming as an acknowledged SoS, crop simulation with DSSAT, and the reason
DSSAT is central to the work presented in this thesis.

\begin{figure}[ht]
  \centering
  \includegraphics[width=0.92\textwidth]{firas/traditional_vs_sos_farming.png}
  \caption{Comparison between traditional reactive farm automation (linear
    sensing--control--actuation chain) and a System-of-Systems smart farm, where
    weather, sensors, and crop simulation feed a shared decision layer that
    commands actuators in a closed loop.}
  \label{fig:traditional_vs_sos}
\end{figure}


\section{System of Systems}
\label{sec:sos}

A System of Systems (SoS) is not simply a large system. It is a set of
independent systems that keep their own purpose and operation, yet work together
to achieve something broader~\cite{maier1998architecting, jamshidi2008systems}.
A weather service, a sensor network, and a crop model can each function without
the others, but when they are combined under a common management goal, they can
support decisions that none of them could support alone~\cite{madni2014integration}.

Agriculture offers a clear example. A farmer may use soil probes from one
company, weather data from a public API, a crop model developed by a research
community, and irrigation equipment installed by a local provider. These tools
were not designed as one package, yet the farmer still needs them to fit
together. If they do not, management becomes fragmented: irrigation reacts to
moisture, fertilisation follows a calendar, and yield expectations ignore the
season's actual weather~\cite{nielsen2015sos, krisnawijaya2024architecture}.


\subsection{SoS Characteristics}

Maier~\cite{maier1998architecting} proposed five characteristics that help
distinguish an SoS from an ordinary integrated system. Table~\ref{tab:sos_characteristics}
lists them and translates each one into an agricultural setting.

\begin{table}[ht]
  \centering
  \caption{System-of-Systems characteristics and their meaning in smart farming.}
  \label{tab:sos_characteristics}
  \begin{tabular}{p{3.2cm} p{4.3cm} p{6.3cm}}
    \toprule
    \textbf{Characteristic} & \textbf{General meaning} & \textbf{Example in agriculture} \\
    \midrule
    Operational independence &
    Each system can function on its own &
    A weather API remains available even if irrigation valves are offline \\[4pt]

    Managerial independence &
    Different authorities own and manage subsystems &
    Sensors, DSSAT, and actuators are developed and maintained by different organisations \\[4pt]

    Geographic distribution &
    Systems are spread across locations &
    Field sensors, cloud services, and farmer terminals are spatially separated \\[4pt]

    Evolutionary development &
    Subsystems can be added, removed, or upgraded independently &
    A drone module or new crop model can be integrated without rebuilding the entire platform \\[4pt]

    Emergent behavior &
    Global behaviour exceeds local subsystem design &
    Yield and resource efficiency emerge from the interaction of weather, soil, crop, and management decisions \\
    \bottomrule
  \end{tabular}
\end{table}

\paragraph{Operational independence.}
Each part of an SoS should remain useful on its own. A weather station keeps
recording even when irrigation is switched off. A crop model can still project
yield from historical weather even if live sensors fail. That independence matters
in agriculture because farms rarely replace all of their equipment at once.

\paragraph{Managerial independence.}
In practice, no single organisation controls every layer of a smart farm. Sensor
manufacturers, weather providers, software platforms, and public crop models
such as DSSAT follow their own development paths. Integration therefore depends
on interfaces and standards, not on one central redesign of the whole
stack~\cite{madni2014integration}.

\paragraph{Geographic distribution.}
Field sensors sit in the soil, gateways may be located at the farm edge, cloud
services run elsewhere, and the farmer may check results on a phone. Coordination
has to happen across distance and networks, not inside one local
machine~\cite{nielsen2015sos}.

\paragraph{Evolutionary development.}
Farm technologies do not age at the same speed. A farmer may upgrade soil probes
one year, adopt a new forecast service the next, and keep the same irrigation
hardware for a decade. An SoS architecture needs to absorb those changes without
forcing a complete restart~\cite{nielsen2015sos, madni2014integration}.

\paragraph{Emergent behavior.}
Yield, water productivity, and nitrogen efficiency are not produced by any single
device. They emerge from the interaction between weather, soil conditions, crop
development, and management over time~\cite{jamshidi2008systems, boote2019cropmodelling}.
That is also why agricultural decision-making is difficult: the outcome farmers
care about appears at the system level, not inside one subsystem.


\subsection{Classification of SoS}
\label{subsec:sos_classification}

Not all Systems of Systems are organised in the same way. Maier~\cite{maier1998architecting}
classified them according to how much central authority exists over the
constituent systems. He originally identified three types---directed,
collaborative, and virtual---and a fourth type, \emph{acknowledged}, was added
later by Dahmann and Baldwin and incorporated into ISO~21839~\cite{nielsen2015sos}.
Table~\ref{tab:sos_taxonomy} gives a short overview.

\begin{table}[ht]
  \centering
  \caption{Classification of Systems of Systems (Maier; extended by Dahmann and Baldwin).}
  \label{tab:sos_taxonomy}
  \begin{tabular}{p{2.5cm} p{4.5cm} p{6.5cm}}
    \toprule
    \textbf{Type} & \textbf{Management structure} & \textbf{Typical agricultural interpretation} \\
    \midrule
    Directed &
    A central authority creates and manages the SoS &
    A fully integrated farm platform owned and operated by one company \\[4pt]

    Acknowledged &
    A recognised SoS manager coordinates independently owned systems &
    A farm-management platform integrates vendor sensors, weather data, DSSAT, and actuators through agreed interfaces \\[4pt]

    Collaborative &
    No central manager; systems cooperate voluntarily &
    A farmer manually combines advisory tools, weather services, and field monitoring applications \\[4pt]

    Virtual &
    No manager and no shared purpose &
    Independent technologies used side by side without real integration \\
    \bottomrule
  \end{tabular}
\end{table}

In a directed SoS, the central authority defines the mission and the subsystems
are built to serve it. In an acknowledged SoS, a coordinating layer exists, but
the subsystems still belong to different owners and evolve on their own;
integration happens through agreement rather than command. Collaborative SoS
depend on voluntary cooperation, while virtual SoS show emergent behaviour even
when the parts were never intended to form a unified
system~\cite{maier1998architecting, nielsen2015sos}.

For modern smart-farming platforms, the acknowledged type is especially
convincing. A farm dashboard, digital twin, or decision-support service can
coordinate tools it does not own. That description fits agriculture well: the
farmer or platform operator manages the overall mission, but the sensors,
models, and actuators remain separate products developed by separate actors.


\section{Smart Farming as an Acknowledged SoS}
\label{sec:smart_farming_sos}

Precision agriculture began with a simple idea: fields are not uniform, so
management should not be uniform either~\cite{aarzoo2025precision, kushwaha2024precision}.
Instead of applying the same irrigation or fertiliser rate everywhere, farmers
can adjust treatment to local soil conditions, moisture levels, crop status, and
weather~\cite{kushwaha2024precision}. Over time, that idea grew into smart
farming: a setting in which sensing, connectivity, data analysis, simulation,
and automated actuation are woven into one management
environment~\cite{hayat2025iot, jadhav2023smartfarming}.

GPS, GIS, remote sensing, and IoT platforms now make it possible to observe
field variability with much greater detail than before~\cite{sharma2020ml,
adewuyi2024datascience}. The result is no longer a single machine or a single
software package, but a network of tools that must cooperate. That is why smart
farming is better described as an SoS than as a standalone automation
system. Table~\ref{tab:traditional_vs_sos} contrasts the older reactive pattern
with this integrated view.

\begin{table}[ht]
  \centering
  \caption{Comparison between traditional farm automation and System-of-Systems smart farming.}
  \label{tab:traditional_vs_sos}
  \begin{tabular}{p{3.5cm} p{5.3cm} p{5.3cm}}
    \toprule
    \textbf{Dimension} & \textbf{Traditional approach} & \textbf{SoS smart farming} \\
    \midrule
    Architecture &
    Isolated chain: sensor $\rightarrow$ controller $\rightarrow$ actuator &
    Multiple independent subsystems connected through a shared decision layer \\[4pt]

    Decision logic &
    Reactive rule-based control &
    Predictive and integrative decision-making \\[4pt]

    Information flow &
    One input, one action &
    Weather, field observations, and crop-model outputs are fused \\[4pt]

    Time horizon &
    Immediate response &
    Short-term control combined with seasonal crop response \\[4pt]

    Subsystem autonomy &
    Low; components designed as one package &
    High; each subsystem can exist independently \\[4pt]

    Scalability &
    Limited to one loop or one field function &
    New sensors, models, or actuators can be added incrementally \\[4pt]

    Closed loop &
    Local physical feedback only &
    Physical feedback plus model-based anticipation \\
    \bottomrule
  \end{tabular}
\end{table}


\subsection{Constituent Systems in Smart Farming}

Even when they are not sold as separate products, the main parts of a smart farm
play distinct roles. Table~\ref{tab:smart_farm_subsystems} summarises the
subsystems that appear repeatedly in the literature and in practical deployments.

\begin{table}[ht]
  \centering
  \caption{Main constituent systems in a smart farming System-of-Systems.}
  \label{tab:smart_farm_subsystems}
  \begin{tabular}{p{2.8cm} p{3.5cm} p{3.5cm} p{3.8cm}}
    \toprule
    \textbf{Subsystem} & \textbf{Main role} & \textbf{Typical inputs} & \textbf{Typical outputs} \\
    \midrule
    Weather system &
    Climate forcing and forecasting &
    Historical weather, numerical forecasts &
    Temperature, rainfall, solar radiation \\[4pt]

    Sensor network &
    Field observation &
    Soil moisture, temperature, nutrient status &
    Ground-truth state measurements \\[4pt]

    Crop simulation model &
    Growth and stress prediction &
    Weather, soil, cultivar, management actions &
    Crop condition, stress indicators, yield projections \\[4pt]

    Decision layer &
    Policy optimisation &
    Observations, model outputs, objectives &
    Irrigation and fertilisation decisions \\[4pt]

    Actuators &
    Execution of management actions &
    Control commands &
    Applied water and nutrients \\
    \bottomrule
  \end{tabular}
\end{table}

\paragraph{Weather system.}
Weather remains the main external driver of crop growth. Local stations can
supply temperature, rainfall, and radiation, while services such as the NASA
POWER API can fill gaps with historical or near-real-time
data~\cite{hayat2025iot, codeluppi2020lorafarm}. Without reliable weather input,
both irrigation planning and crop modelling become unreliable.

\paragraph{Sensor network.}
Soil sensors based on capacitance or time-domain reflectometry (TDR) can track
moisture, temperature, and nitrogen status at different
depths~\cite{hayat2025iot, jadhav2023smartfarming}. They provide the field-level
observations that make site-specific management possible.

\paragraph{Communication infrastructure.}
Distributed farms need a way to move data from the field to the decision layer.
LPWAN technologies, especially LoRaWAN, are often used because they cover long
distances with low energy use~\cite{codeluppi2020lorafarm, jadhav2023smartfarming}.
In an SoS view, this communication layer is what turns separate devices into a
connected system.

\paragraph{Crop simulation model.}
Sensors show what is happening now; crop models help estimate what management
will mean later in the season. They connect weather, soil, and farmer actions to
likely crop development, stress, and harvest
outcomes~\cite{boote2019cropmodelling, hoogenboom2019dssat}. They extend the
decision horizon beyond the latest sensor reading.

\paragraph{Decision layer.}
At some point, the farm must choose an action: irrigate or wait, fertilise or
hold. That choice may come from a rule engine, an optimisation tool, a human
adviser, or---as in this thesis---a learning-based controller. Whatever form it
takes, the decision layer is where the SoS becomes operational.

\paragraph{Actuators.}
Irrigation valves and variable-rate fertiliser applicators turn decisions into
physical changes in the field~\cite{hayat2025iot}. Once water or nitrogen is
applied, the crop and soil state changes, and the loop starts again.


\subsection{Why Smart Farming Fits the Acknowledged SoS Model}

Smart farming resembles an acknowledged SoS for a simple reason: coordination
exists, but ownership stays divided. A platform or farm manager integrates
weather, sensors, models, and actuators, yet those components usually come from
different suppliers and remain replaceable. A farmer can change sensor brand or
forecast provider without rebuilding the entire system.

That does not mean every farm already works this way. Many operations still look
more like the traditional chain in Figure~\ref{fig:traditional_vs_sos}. The
acknowledged SoS model is therefore best seen as the direction in which mature
smart-farming platforms are heading: separate tools, shared
mission~\cite{krisnawijaya2024architecture, nielsen2015sos}.

There is another practical complication. The data arriving from these subsystems
do not share the same form. Sensor streams are time series, drone imagery is
visual and high-dimensional, and crop-model outputs may appear as state
variables or structured reports~\cite{krisnawijaya2024architecture,
adewuyi2024datascience}. A workable SoS therefore needs both communication
infrastructure and a decision layer able to bring those different sources
together.


\section{Crop Simulation and DSSAT}
\label{sec:dssat}

Sensors alone cannot answer every management question. A moisture reading may
show that the soil is dry today, but it does not by itself predict final yield,
nitrogen losses, or crop stress two weeks later. Crop simulation models fill
that gap by representing how plants, soil water, and nutrients respond over
time~\cite{boote2019cropmodelling}.

Among the available platforms, the Decision Support System for Agrotechnology
Transfer (DSSAT) is one of the most established. It has been used for decades to
combine variety, environment, and management information in dynamic crop
simulations~\cite{hoogenboom2019dssat, sachin2019dssat}. Its main value is not
a single yield number at harvest, but the ability to follow how daily weather
and management actions shape crop response over the season.


\subsection{The G$\times$E$\times$M Framework}

DSSAT is built around the Genotype--Environment--Management (G$\times$E$\times$M)
framework. Crop performance depends on all three. A high-yielding cultivar still
needs suitable weather and reasonable management; good irrigation cannot fully
compensate for a poor variety or badly timed
operations~\cite{boote2019cropmodelling}. Table~\ref{tab:gxem} breaks down the
three components.

\begin{table}[ht]
  \centering
  \caption{The Genotype--Environment--Management (G$\times$E$\times$M) framework in DSSAT.}
  \label{tab:gxem}
  \begin{tabular}{p{2.5cm} p{4.3cm} p{6.5cm}}
    \toprule
    \textbf{Factor} & \textbf{Meaning} & \textbf{Examples in DSSAT} \\
    \midrule
    Genotype (G) &
    Crop genetic potential and development pattern &
    Cultivar parameters, phenology, grain-fill capacity \\[4pt]

    Environment (E) &
    Climatic and edaphic conditions &
    Daily weather, soil layers, initial water and nitrogen status \\[4pt]

    Management (M) &
    Farmer interventions during the season &
    Planting date, irrigation depth, nitrogen applications \\
    \bottomrule
  \end{tabular}
\end{table}

This matters for decision support because management actions do not influence
yield directly. Irrigation changes soil moisture; fertilisation changes nutrient
availability; those changes then affect how the crop grows, whether it becomes
stressed, and what is harvested at the end of the season. DSSAT makes that chain
explicit.


\subsection{The DSSAT Cropping System Model}

At the heart of DSSAT is the Cropping System Model, usually referred to as CSM.
Instead of jumping straight to a harvest number, CSM moves through the season one
day at a time. On each day, it reads the weather, updates soil moisture and
nutrient conditions, applies any field operations such as irrigation or
fertiliser, and then updates the state of the
crop~\cite{hoogenboom2019dssat}.

For the maize experiments in this thesis, the relevant crop model is
CERES-Maize. In everyday terms, it tracks how the crop develops after planting,
how well it copes with dry or nutrient-poor conditions, and how much grain is
likely to be produced by harvest. That makes it useful for the kind of management
questions studied here, where water and nitrogen decisions need to be judged not
only for their immediate effect, but also for their impact over the rest of the
season~\cite{boote2019cropmodelling, sachin2019dssat}.

This is also what distinguishes DSSAT from a simple prediction formula. A basic
model might estimate yield from a few inputs, but it would miss the interactions
that make farm management difficult in practice: too little water early in the
season, too much nitrogen before rain, or the opposite. CSM is built to represent
those interactions rather than ignore them.

In most published DSSAT studies, the management plan is decided in advance and the
model is run as a full-season experiment. In this thesis, the same crop-response
logic is used, but irrigation and nitrogen applications may change over time. The
software details of that setup are presented in the methodology chapters; here,
the important point is simply that CSM provides the crop-response core around
which the rest of the system is built.


\section{Why DSSAT for This Work}
\label{sec:dssat_architecture}

Look again at Figure~\ref{fig:traditional_vs_sos}. Weather tells the system what
the climate is doing. Sensors describe the current field state. DSSAT connects
those conditions to likely crop response under different irrigation and
fertilisation choices. It is the part of the SoS that translates management
decisions into crop outcomes.

That is the question at the centre of this thesis. The work is not about
building a complete smart farm with its own communication stack, drone service,
and farm dashboard. It is about what happens inside the SoS when yield, water
use, and nitrogen efficiency cannot all be maximised at once: how should
irrigation and fertilisation be adjusted over the season? DSSAT makes that
question concrete because it allows those trade-offs to be explored against a
realistic crop-response model~\cite{hoogenboom2019dssat, sachin2019dssat}.

The next chapter introduces the reinforcement-learning framework used to study
that problem. The methodology chapters then explain how daily management
decisions are exchanged with the simulator in practice.


\section{Conclusion}
\label{sec:ch1_conclusion}

Modern farming is moving away from isolated devices that react to a single
reading and toward environments in which weather services, sensors, crop models,
and actuators have to work together. This chapter used the System-of-Systems view
to describe that shift and to explain why coordination alone is not enough:
some part of the system still has to decide how water, nitrogen, and crop
outcomes should be balanced over a season.

Smart farming fits the acknowledged SoS model because farms increasingly rely on
tools they do not fully own, yet still need those tools to serve a common
management goal. DSSAT enters the picture as the crop-simulation part of that
picture. It does not replace sensors or actuators, but it helps answer what
happens to the crop when management changes under real weather and soil
conditions.

That is why DSSAT matters for this thesis. The work is less about assembling a
complete commercial smart farm and more about studying a difficult decision
problem inside that broader setting: how irrigation and fertilisation should
adapt when yield, water use, and nitrogen efficiency pull in different
directions. The next chapter takes up that decision problem directly and
introduces the reinforcement-learning framework used to address it.

% Do NOT compile this file alone. Set main.tex as the Overleaf main document.
% Required packages in main.tex: booktabs, graphicx, array
% Required in preamble: \graphicspath{{firas/}{figures/}}
% =============================================================================

\chapter{Background and Context}
\label{ch:background}

On many farms today, automation still works in simple chains. A sensor reads
soil moisture, a controller checks a threshold, and a pump switches on. Each
step makes sense on its own, but the loop rarely looks beyond the immediate
reading. It does not account for tomorrow's weather, the crop's growth stage, or
the longer-term trade-off between saving water and protecting yield.

That limitation becomes more visible as farms adopt more technologies. A single
field may now depend on soil sensors, weather data, communication networks,
simulation tools, and automated irrigation hardware, often supplied by different
vendors and updated on different schedules. The challenge is no longer to operate
one device well, but to make many independent parts work toward the same goal:
producing food efficiently while using water and nutrients responsibly. This is
where the \emph{System-of-Systems} (SoS) perspective becomes useful.

Figure~\ref{fig:traditional_vs_sos} illustrates the difference. The upper part
of the figure shows the traditional pattern: one input, one decision, one action.
The lower part shows a smart-farming SoS, in which weather data, field
observations, and crop-growth prediction feed a shared decision layer that
commands actuators and then receives updated feedback from the field. The rest
of this chapter builds the background needed for that view: SoS theory, smart
farming as an acknowledged SoS, crop simulation with DSSAT, and the reason
DSSAT is central to the work presented in this thesis.

\begin{figure}[ht]
  \centering
  \includegraphics[width=0.92\textwidth]{firas/traditional_vs_sos_farming.png}
  \caption{Comparison between traditional reactive farm automation (linear
    sensing--control--actuation chain) and a System-of-Systems smart farm, where
    weather, sensors, and crop simulation feed a shared decision layer that
    commands actuators in a closed loop.}
  \label{fig:traditional_vs_sos}
\end{figure}


\section{System of Systems}
\label{sec:sos}

A System of Systems (SoS) is not simply a large system. It is a set of
independent systems that keep their own purpose and operation, yet work together
to achieve something broader~\cite{maier1998architecting, jamshidi2008systems}.
A weather service, a sensor network, and a crop model can each function without
the others, but when they are combined under a common management goal, they can
support decisions that none of them could support alone~\cite{madni2014integration}.

Agriculture offers a clear example. A farmer may use soil probes from one
company, weather data from a public API, a crop model developed by a research
community, and irrigation equipment installed by a local provider. These tools
were not designed as one package, yet the farmer still needs them to fit
together. If they do not, management becomes fragmented: irrigation reacts to
moisture, fertilisation follows a calendar, and yield expectations ignore the
season's actual weather~\cite{nielsen2015sos, krisnawijaya2024architecture}.


\subsection{SoS Characteristics}

Maier~\cite{maier1998architecting} proposed five characteristics that help
distinguish an SoS from an ordinary integrated system. Table~\ref{tab:sos_characteristics}
lists them and translates each one into an agricultural setting.

\begin{table}[ht]
  \centering
  \caption{System-of-Systems characteristics and their meaning in smart farming.}
  \label{tab:sos_characteristics}
  \begin{tabular}{p{3.2cm} p{4.3cm} p{6.3cm}}
    \toprule
    \textbf{Characteristic} & \textbf{General meaning} & \textbf{Example in agriculture} \\
    \midrule
    Operational independence &
    Each system can function on its own &
    A weather API remains available even if irrigation valves are offline \\[4pt]

    Managerial independence &
    Different authorities own and manage subsystems &
    Sensors, DSSAT, and actuators are developed and maintained by different organisations \\[4pt]

    Geographic distribution &
    Systems are spread across locations &
    Field sensors, cloud services, and farmer terminals are spatially separated \\[4pt]

    Evolutionary development &
    Subsystems can be added, removed, or upgraded independently &
    A drone module or new crop model can be integrated without rebuilding the entire platform \\[4pt]

    Emergent behavior &
    Global behaviour exceeds local subsystem design &
    Yield and resource efficiency emerge from the interaction of weather, soil, crop, and management decisions \\
    \bottomrule
  \end{tabular}
\end{table}

\paragraph{Operational independence.}
Each part of an SoS should remain useful on its own. A weather station keeps
recording even when irrigation is switched off. A crop model can still project
yield from historical weather even if live sensors fail. That independence matters
in agriculture because farms rarely replace all of their equipment at once.

\paragraph{Managerial independence.}
In practice, no single organisation controls every layer of a smart farm. Sensor
manufacturers, weather providers, software platforms, and public crop models
such as DSSAT follow their own development paths. Integration therefore depends
on interfaces and standards, not on one central redesign of the whole
stack~\cite{madni2014integration}.

\paragraph{Geographic distribution.}
Field sensors sit in the soil, gateways may be located at the farm edge, cloud
services run elsewhere, and the farmer may check results on a phone. Coordination
has to happen across distance and networks, not inside one local
machine~\cite{nielsen2015sos}.

\paragraph{Evolutionary development.}
Farm technologies do not age at the same speed. A farmer may upgrade soil probes
one year, adopt a new forecast service the next, and keep the same irrigation
hardware for a decade. An SoS architecture needs to absorb those changes without
forcing a complete restart~\cite{nielsen2015sos, madni2014integration}.

\paragraph{Emergent behavior.}
Yield, water productivity, and nitrogen efficiency are not produced by any single
device. They emerge from the interaction between weather, soil conditions, crop
development, and management over time~\cite{jamshidi2008systems, boote2019cropmodelling}.
That is also why agricultural decision-making is difficult: the outcome farmers
care about appears at the system level, not inside one subsystem.


\subsection{Classification of SoS}
\label{subsec:sos_classification}

Not all Systems of Systems are organised in the same way. Maier~\cite{maier1998architecting}
classified them according to how much central authority exists over the
constituent systems. He originally identified three types---directed,
collaborative, and virtual---and a fourth type, \emph{acknowledged}, was added
later by Dahmann and Baldwin and incorporated into ISO~21839~\cite{nielsen2015sos}.
Table~\ref{tab:sos_taxonomy} gives a short overview.

\begin{table}[ht]
  \centering
  \caption{Classification of Systems of Systems (Maier; extended by Dahmann and Baldwin).}
  \label{tab:sos_taxonomy}
  \begin{tabular}{p{2.5cm} p{4.5cm} p{6.5cm}}
    \toprule
    \textbf{Type} & \textbf{Management structure} & \textbf{Typical agricultural interpretation} \\
    \midrule
    Directed &
    A central authority creates and manages the SoS &
    A fully integrated farm platform owned and operated by one company \\[4pt]

    Acknowledged &
    A recognised SoS manager coordinates independently owned systems &
    A farm-management platform integrates vendor sensors, weather data, DSSAT, and actuators through agreed interfaces \\[4pt]

    Collaborative &
    No central manager; systems cooperate voluntarily &
    A farmer manually combines advisory tools, weather services, and field monitoring applications \\[4pt]

    Virtual &
    No manager and no shared purpose &
    Independent technologies used side by side without real integration \\
    \bottomrule
  \end{tabular}
\end{table}

In a directed SoS, the central authority defines the mission and the subsystems
are built to serve it. In an acknowledged SoS, a coordinating layer exists, but
the subsystems still belong to different owners and evolve on their own;
integration happens through agreement rather than command. Collaborative SoS
depend on voluntary cooperation, while virtual SoS show emergent behaviour even
when the parts were never intended to form a unified
system~\cite{maier1998architecting, nielsen2015sos}.

For modern smart-farming platforms, the acknowledged type is especially
convincing. A farm dashboard, digital twin, or decision-support service can
coordinate tools it does not own. That description fits agriculture well: the
farmer or platform operator manages the overall mission, but the sensors,
models, and actuators remain separate products developed by separate actors.


\section{Smart Farming as an Acknowledged SoS}
\label{sec:smart_farming_sos}

Precision agriculture began with a simple idea: fields are not uniform, so
management should not be uniform either~\cite{aarzoo2025precision, kushwaha2024precision}.
Instead of applying the same irrigation or fertiliser rate everywhere, farmers
can adjust treatment to local soil conditions, moisture levels, crop status, and
weather~\cite{kushwaha2024precision}. Over time, that idea grew into smart
farming: a setting in which sensing, connectivity, data analysis, simulation,
and automated actuation are woven into one management
environment~\cite{hayat2025iot, jadhav2023smartfarming}.

GPS, GIS, remote sensing, and IoT platforms now make it possible to observe
field variability with much greater detail than before~\cite{sharma2020ml,
adewuyi2024datascience}. The result is no longer a single machine or a single
software package, but a network of tools that must cooperate. That is why smart
farming is better described as an SoS than as a standalone automation
system. Table~\ref{tab:traditional_vs_sos} contrasts the older reactive pattern
with this integrated view.

\begin{table}[ht]
  \centering
  \caption{Comparison between traditional farm automation and System-of-Systems smart farming.}
  \label{tab:traditional_vs_sos}
  \begin{tabular}{p{3.5cm} p{5.3cm} p{5.3cm}}
    \toprule
    \textbf{Dimension} & \textbf{Traditional approach} & \textbf{SoS smart farming} \\
    \midrule
    Architecture &
    Isolated chain: sensor $\rightarrow$ controller $\rightarrow$ actuator &
    Multiple independent subsystems connected through a shared decision layer \\[4pt]

    Decision logic &
    Reactive rule-based control &
    Predictive and integrative decision-making \\[4pt]

    Information flow &
    One input, one action &
    Weather, field observations, and crop-model outputs are fused \\[4pt]

    Time horizon &
    Immediate response &
    Short-term control combined with seasonal crop response \\[4pt]

    Subsystem autonomy &
    Low; components designed as one package &
    High; each subsystem can exist independently \\[4pt]

    Scalability &
    Limited to one loop or one field function &
    New sensors, models, or actuators can be added incrementally \\[4pt]

    Closed loop &
    Local physical feedback only &
    Physical feedback plus model-based anticipation \\
    \bottomrule
  \end{tabular}
\end{table}


\subsection{Constituent Systems in Smart Farming}

Even when they are not sold as separate products, the main parts of a smart farm
play distinct roles. Table~\ref{tab:smart_farm_subsystems} summarises the
subsystems that appear repeatedly in the literature and in practical deployments.

\begin{table}[ht]
  \centering
  \caption{Main constituent systems in a smart farming System-of-Systems.}
  \label{tab:smart_farm_subsystems}
  \begin{tabular}{p{2.8cm} p{3.5cm} p{3.5cm} p{3.8cm}}
    \toprule
    \textbf{Subsystem} & \textbf{Main role} & \textbf{Typical inputs} & \textbf{Typical outputs} \\
    \midrule
    Weather system &
    Climate forcing and forecasting &
    Historical weather, numerical forecasts &
    Temperature, rainfall, solar radiation \\[4pt]

    Sensor network &
    Field observation &
    Soil moisture, temperature, nutrient status &
    Ground-truth state measurements \\[4pt]

    Crop simulation model &
    Growth and stress prediction &
    Weather, soil, cultivar, management actions &
    Crop condition, stress indicators, yield projections \\[4pt]

    Decision layer &
    Policy optimisation &
    Observations, model outputs, objectives &
    Irrigation and fertilisation decisions \\[4pt]

    Actuators &
    Execution of management actions &
    Control commands &
    Applied water and nutrients \\
    \bottomrule
  \end{tabular}
\end{table}

\paragraph{Weather system.}
Weather remains the main external driver of crop growth. Local stations can
supply temperature, rainfall, and radiation, while services such as the NASA
POWER API can fill gaps with historical or near-real-time
data~\cite{hayat2025iot, codeluppi2020lorafarm}. Without reliable weather input,
both irrigation planning and crop modelling become unreliable.

\paragraph{Sensor network.}
Soil sensors based on capacitance or time-domain reflectometry (TDR) can track
moisture, temperature, and nitrogen status at different
depths~\cite{hayat2025iot, jadhav2023smartfarming}. They provide the field-level
observations that make site-specific management possible.

\paragraph{Communication infrastructure.}
Distributed farms need a way to move data from the field to the decision layer.
LPWAN technologies, especially LoRaWAN, are often used because they cover long
distances with low energy use~\cite{codeluppi2020lorafarm, jadhav2023smartfarming}.
In an SoS view, this communication layer is what turns separate devices into a
connected system.

\paragraph{Crop simulation model.}
Sensors show what is happening now; crop models help estimate what management
will mean later in the season. They connect weather, soil, and farmer actions to
likely crop development, stress, and harvest
outcomes~\cite{boote2019cropmodelling, hoogenboom2019dssat}. They extend the
decision horizon beyond the latest sensor reading.

\paragraph{Decision layer.}
At some point, the farm must choose an action: irrigate or wait, fertilise or
hold. That choice may come from a rule engine, an optimisation tool, a human
adviser, or---as in this thesis---a learning-based controller. Whatever form it
takes, the decision layer is where the SoS becomes operational.

\paragraph{Actuators.}
Irrigation valves and variable-rate fertiliser applicators turn decisions into
physical changes in the field~\cite{hayat2025iot}. Once water or nitrogen is
applied, the crop and soil state changes, and the loop starts again.


\subsection{Why Smart Farming Fits the Acknowledged SoS Model}

Smart farming resembles an acknowledged SoS for a simple reason: coordination
exists, but ownership stays divided. A platform or farm manager integrates
weather, sensors, models, and actuators, yet those components usually come from
different suppliers and remain replaceable. A farmer can change sensor brand or
forecast provider without rebuilding the entire system.

That does not mean every farm already works this way. Many operations still look
more like the traditional chain in Figure~\ref{fig:traditional_vs_sos}. The
acknowledged SoS model is therefore best seen as the direction in which mature
smart-farming platforms are heading: separate tools, shared
mission~\cite{krisnawijaya2024architecture, nielsen2015sos}.

There is another practical complication. The data arriving from these subsystems
do not share the same form. Sensor streams are time series, drone imagery is
visual and high-dimensional, and crop-model outputs may appear as state
variables or structured reports~\cite{krisnawijaya2024architecture,
adewuyi2024datascience}. A workable SoS therefore needs both communication
infrastructure and a decision layer able to bring those different sources
together.


\section{Crop Simulation and DSSAT}
\label{sec:dssat}

Sensors alone cannot answer every management question. A moisture reading may
show that the soil is dry today, but it does not by itself predict final yield,
nitrogen losses, or crop stress two weeks later. Crop simulation models fill
that gap by representing how plants, soil water, and nutrients respond over
time~\cite{boote2019cropmodelling}.

Among the available platforms, the Decision Support System for Agrotechnology
Transfer (DSSAT) is one of the most established. It has been used for decades to
combine variety, environment, and management information in dynamic crop
simulations~\cite{hoogenboom2019dssat, sachin2019dssat}. Its main value is not
a single yield number at harvest, but the ability to follow how daily weather
and management actions shape crop response over the season.


\subsection{The G$\times$E$\times$M Framework}

DSSAT is built around the Genotype--Environment--Management (G$\times$E$\times$M)
framework. Crop performance depends on all three. A high-yielding cultivar still
needs suitable weather and reasonable management; good irrigation cannot fully
compensate for a poor variety or badly timed
operations~\cite{boote2019cropmodelling}. Table~\ref{tab:gxem} breaks down the
three components.

\begin{table}[ht]
  \centering
  \caption{The Genotype--Environment--Management (G$\times$E$\times$M) framework in DSSAT.}
  \label{tab:gxem}
  \begin{tabular}{p{2.5cm} p{4.3cm} p{6.5cm}}
    \toprule
    \textbf{Factor} & \textbf{Meaning} & \textbf{Examples in DSSAT} \\
    \midrule
    Genotype (G) &
    Crop genetic potential and development pattern &
    Cultivar parameters, phenology, grain-fill capacity \\[4pt]

    Environment (E) &
    Climatic and edaphic conditions &
    Daily weather, soil layers, initial water and nitrogen status \\[4pt]

    Management (M) &
    Farmer interventions during the season &
    Planting date, irrigation depth, nitrogen applications \\
    \bottomrule
  \end{tabular}
\end{table}

This matters for decision support because management actions do not influence
yield directly. Irrigation changes soil moisture; fertilisation changes nutrient
availability; those changes then affect how the crop grows, whether it becomes
stressed, and what is harvested at the end of the season. DSSAT makes that chain
explicit.


\subsection{The DSSAT Cropping System Model}

At the heart of DSSAT is the Cropping System Model, usually referred to as CSM.
Instead of jumping straight to a harvest number, CSM moves through the season one
day at a time. On each day, it reads the weather, updates soil moisture and
nutrient conditions, applies any field operations such as irrigation or
fertiliser, and then updates the state of the
crop~\cite{hoogenboom2019dssat}.

For the maize experiments in this thesis, the relevant crop model is
CERES-Maize. In everyday terms, it tracks how the crop develops after planting,
how well it copes with dry or nutrient-poor conditions, and how much grain is
likely to be produced by harvest. That makes it useful for the kind of management
questions studied here, where water and nitrogen decisions need to be judged not
only for their immediate effect, but also for their impact over the rest of the
season~\cite{boote2019cropmodelling, sachin2019dssat}.

This is also what distinguishes DSSAT from a simple prediction formula. A basic
model might estimate yield from a few inputs, but it would miss the interactions
that make farm management difficult in practice: too little water early in the
season, too much nitrogen before rain, or the opposite. CSM is built to represent
those interactions rather than ignore them.

In most published DSSAT studies, the management plan is decided in advance and the
model is run as a full-season experiment. In this thesis, the same crop-response
logic is used, but irrigation and nitrogen applications may change over time. The
software details of that setup are presented in the methodology chapters; here,
the important point is simply that CSM provides the crop-response core around
which the rest of the system is built.


\section{Why DSSAT for This Work}
\label{sec:dssat_architecture}

Look again at Figure~\ref{fig:traditional_vs_sos}. Weather tells the system what
the climate is doing. Sensors describe the current field state. DSSAT connects
those conditions to likely crop response under different irrigation and
fertilisation choices. It is the part of the SoS that translates management
decisions into crop outcomes.

That is the question at the centre of this thesis. The work is not about
building a complete smart farm with its own communication stack, drone service,
and farm dashboard. It is about what happens inside the SoS when yield, water
use, and nitrogen efficiency cannot all be maximised at once: how should
irrigation and fertilisation be adjusted over the season? DSSAT makes that
question concrete because it allows those trade-offs to be explored against a
realistic crop-response model~\cite{hoogenboom2019dssat, sachin2019dssat}.

The next chapter introduces the reinforcement-learning framework used to study
that problem. The methodology chapters then explain how daily management
decisions are exchanged with the simulator in practice.


\section{Conclusion}
\label{sec:ch1_conclusion}

Modern farming is moving away from isolated devices that react to a single
reading and toward environments in which weather services, sensors, crop models,
and actuators have to work together. This chapter used the System-of-Systems view
to describe that shift and to explain why coordination alone is not enough:
some part of the system still has to decide how water, nitrogen, and crop
outcomes should be balanced over a season.

Smart farming fits the acknowledged SoS model because farms increasingly rely on
tools they do not fully own, yet still need those tools to serve a common
management goal. DSSAT enters the picture as the crop-simulation part of that
picture. It does not replace sensors or actuators, but it helps answer what
happens to the crop when management changes under real weather and soil
conditions.

That is why DSSAT matters for this thesis. The work is less about assembling a
complete commercial smart farm and more about studying a difficult decision
problem inside that broader setting: how irrigation and fertilisation should
adapt when yield, water use, and nitrogen efficiency pull in different
directions. The next chapter takes up that decision problem directly and
introduces the reinforcement-learning framework used to address it.
